\documentclass{book} % Definisi jenis dokumen

%%%%% Definisi paket-paket yang digunakan %%%%%
\usepackage[T1]{fontenc} % paket encoding huruf (wajib)
\usepackage[utf8]{inputenc} % paket encoding input (wajib)
\usepackage[yyyymmdd,hhmmss]{datetime} % paket tanggal-waktu
\usepackage{graphicx} % paket grafik/gambar
\usepackage[english]{babel} % paket modifikasi label/caption pada bahasa tertentu
\usepackage{geometry} % paket ukuran kertas dan margin
\usepackage{hyperref} % paket referensi url dan hyperref
\usepackage{minted} % paket lingkungan kode sumber
\usepackage{xcolor} % paket definisi warna

%%%%% Pengaturan ukuran kertas dan margin %%%%%
\geometry{
	a4paper,
	left=10mm,
	right=10mm,
	top=15mm,
	bottom=15mm,
}

%%%%% Pengaturan perintah informasi perangkat lunak %%%%%
\newcommand{\ShowOsVersion}{
	\immediate\write18{\unexpanded{foo=`uname -sro` && echo "${foo}" > tmp.tex}}
	\input{tmp}\immediate\write18{rm tmp.tex}
}

\newcommand{\ShowTexVersion}{
	\immediate\write18{\unexpanded{foo=`pdflatex -version | head -n1 | cut -d' ' -f1,2` && echo "${foo}" > tmp.tex}}
	\input{tmp}\immediate\write18{rm tmp.tex}
}

%%%%% Mengganti label "Contents" ke "Daftar Isi" %%%%%
\addto\captionsenglish{\renewcommand{\contentsname}{Daftar Isi}}

%%%%% Pengaturan teks link url/hyperref %%%%%
\hypersetup{
	colorlinks=true, %set true if you want colored links
	linktoc=all,     %set to all if you want both sections and subsections linked
	linkcolor=blue,  %choose some color if you want links to stand out
	urlcolor=blue,   %url color
}

%%%%% Definisi warna baru %%%%%
\definecolor{LightGray}{gray}{0.95}

\begin{document}

    %%%%%%%%%%%%%%%%%%%%%%%%%%%%%%%%%%%%%%%%%%%%%%%%%%%%%%%%%%%%%%%%%

    \frontmatter % untuk halaman cover

    \begin{titlepage}
        \centering % untuk membuat tengah teks

        {
            \LARGE % pakai font besar
            \bf % pakai font BOLD
            Pengenalan LaTeX untuk Pemula
        }

        {\Large \bf Achmadi ST MT}
        \vfill % menambahkan ruang kosong vertikal

		\includegraphics[width=500pt]{images/leafcover}
        \vfill

        \raggedright
        \noindent Buku ini ditulis dengan:\\ % tanda \\ menambahkan garis baru
		OS : \ShowOsVersion \\
		TeX : \ShowTexVersion \\
		Update: {\today} at \currenttime\\
    \end{titlepage}

    %%%%%%%%%%%%%%%%%%%%%%%%%%%%%%%%%%%%%%%%%%%%%%%%%%%%%%%%%%%%%%%%%

	\newpage % halaman baru

	\tableofcontents % daftar isi

	%%%%%%%%%%%%%%%%%%%%%%%%%%%%%%%%%%%%%%%%%%%%%%%%%%%%%%%%%%%%%%%%%

    \newpage
    \chapter{Disclaimer} % memulai chapter baru

    LaTeX adalah proyek typesetting yang dimulai oleh \href{https://en.wikipedia.org/wiki/Leslie_Lamport}{Leslie Lamport}
    sekitar 40 tahun lalu dan sekarang dilanjutkan oleh \href{https://www.latex-project.org/}{latex-project.org}.

    \bigskip % baris kosong pemisah paragraf

    Buku tutorial ini disarikan dari \href{https://latex-tutorial.com/tutorials}{latex-tutorial.com}
    dan \href{https://www.overleaf.com/learn/latex/Tutorials}{overleaf.com} dengan alih bahasa dan penambahan konten sesuai selera penulis.

    %%%%%%%%%%%%%%%%%%%%%%%%%%%%%%%%%%%%%%%%%%%%%%%%%%%%%%%%%%%%%%%%%

	\newpage
	\chapter{Penggunaan Buku}

	\section{Umum} % memulai section baru
	Buku ini dibuat dengan tujuan penggunaan utama sebagai panduan digital untuk mempermudah search dan copy-paste.
	Anda tidak perlu mencetak buku ini ke bentuk kertas.
	Seluruh navigasi buku ini diharapkan menggunakan klik ke hyperlink di Daftar Isi,
	atau menggunakan tampilan \textbf{Index} yang tersedia di \textbf{SideBar} program pembaca PDF yang anda gunakan.

	\section{Petunjuk}
	Beberapa petunjuk yang digunakan di buku ini:
	\begin{itemize} % memulai lingkungan daftar tanpa angka (bullet list)
		\item \textbf{Cetak Tebal}: Menginformasikan identifier (keyword, nama file, dst) yang berada di suatu paragraf
		\item \textit{Cetak Miring}: Bersama simbol panah (->) dan simbol lain, menginformasikan langkah-langkah klik menu/tombol.
		\item \textbf{TIPS:} Menginformasikan hal-hal yang dapat membantu atau pengetahuan tambahan.
		\item \textbf{PERINGATAN:} Menginformasikan hal-hal yang bener-benar harus diperhatikan.
	\end{itemize}

	\section{Penulisan Kode}
	Untuk penulisan kode, akan digunakan sebagai berikut:

    % lingkungan penulisan kode sumber
    \begin{minted}[frame=lines,framesep=2mm,fontsize=\normalsize,bgcolor=LightGray]{latex}
\begin{enumerate}
    \item Pertama
    \item Kedua
\end{enumerate }
    \end{minted}

    \section{Kode Sumber Dokumen}

    Seluruh isi buku ini juga ditulis dengan LaTeX, sehingga jika dibutuhkan contoh, dapat melihat
    kode sumber dokumen ini yang tersedia di \href{https://github.com/mekatronik-achmadi/my_latexbook/blob/master/Modul/LaTex/latex_beginner.tex}{GitHub}.

    \bigskip

    Selain itu, website seperti \href{https://www.overleaf.com/}{Overleaf.com} menyediakan beragam tutorial dan contoh siap pakai.

    %%%%%%%%%%%%%%%%%%%%%%%%%%%%%%%%%%%%%%%%%%%%%%%%%%%%%%%%%%%%%%%%%

    \newpage

    \mainmatter % pindah format halaman dari romawi ke angka (konten utama)
    \chapter{Introduksi}

    \section{Apa itu LaTeX}

    LaTeX adalah bahasa markup (seperti HTML dan Markdown) untuk typeset (mengatur penulisan) dokumen.
    Format berkas dokumen LaTeX umumnya berupa berkas teks yang disimpan di suatu folder
    bersamaan dengan berkas pendukung lain (seperti gambar, pdf, dan sebagainya).
    Berkas teks ini kemudian dikompilasi menjadi berkas akhir (umumnya PDF atau DVI).

    LaTeX sangat populer digunakan untuk antara lain:
    \begin{itemize}
        \item Jurnal, paper, atau artikel akademik.
        \item Laporan Teknis.
        \item Buku atau Novel.
        \item Presentasi.
    \end{itemize}

    \section{Contoh kode LaTeX}

    Berikut contoh kode LaTeX untuk menulis tabel:

    \begin{minted}[frame=lines,framesep=2mm,fontsize=\normalsize,bgcolor=LightGray]{latex}
\begin{center}
    \begin{tabular}{ |c|c|c| }
        \hline
        cell1 & cell2 & cell3 \\
        \hline
        cell4 & cell5 & cell6 \\
        \hline
        cell7 & cell8 & cell9 \\
        \hline
    \end{tabular}
\end{center}
    \end{minted}

    Hasil kode di atas akan menghasilkan sebagai berikut:
    \begin{center}
        \begin{tabular}{ |c|c|c| }
            \hline
            cell1 & cell2 & cell3 \\
            \hline
            cell4 & cell5 & cell6 \\
            \hline
            cell7 & cell8 & cell9 \\
            \hline
        \end{tabular}
    \end{center}

    \section{Komentar dalam LaTeX}

    Selayaknya bahasa pemrograman pada umumnya, LaTeX mengenal mekanisme komentar yang tidak akan dikompilasi ke dokumen akhir.
    Komentar dalam latex adalah komentar per baris yang dimulai dengan karakter persen '\%'.

    Sebagai contoh:

    \begin{minted}[frame=lines,framesep=2mm,fontsize=\normalsize,bgcolor=LightGray]{latex}
\begin{itemize}
    \item First % Ini pertama
    \item Second % Ini kedua
\end{itemize}
    \end{minted}

    Yang akan menghasikan (tanpa ada konten setelah tanda persen) :

    \begin{itemize}
        \item First % Ini pertama
        \item Second % Ini kedua
    \end{itemize}

    Jika dibutuhkan menulis simbol persen, ditambahkan simbol backslash ('\textbackslash').
    Contohnya:

    \begin{minted}[frame=lines,framesep=2mm,fontsize=\normalsize,bgcolor=LightGray]{latex}
Uang senilai 50\% dari Rp.1000 adalah Rp.500
    \end{minted}

    Menghasilkan: \\
    Uang senilai 50\% dari Rp.1000 adalah Rp.500

    \section{Apakah LaTeX termasuk Word-Processor}

    LaTeX sama sekali \textbf{BUKAN} word-processor seperti MS-Word, LO-Writer, OO-Document, Google-Doc, atau yang serupa.
    Perbedaan mendasar adalah:
    \begin{itemize}
        \item Word-Processor menganut paradigma WYSIWYG,
        sedangkan LaTeX adalah markup dan typesetting yang perlu dikompilasi terlebih dahulu.

        \item Word-Processor modern membutuhkan perangkat layout seperti mouse/touchpad,
        sedangkan LaTeX murni programmatikal teks.
    \end{itemize}

    \bigskip

    Keuntungan Word-processor dibandingkan LaTeX:
    \begin{itemize}
        \item Dokumen workflow yang telah dikenalkan sejak kita SD.
        \item Penguna lebih banyak dan lebih umum.
    \end{itemize}

    \bigskip

    Keuntungan LaTeX dibandingkan Word-Processor:
    \begin{itemize}
        \item Satu template dapat digunakan berkali-kali.
        \item Tampilan dokumen diatur template, sehingga penulis bisa fokus pada konten saja.
        \item Hasil penulisan format rumus kompleks bagus.
        \item Hasil penulisan kode programming bagus.
        \item Hasil penulisan bibliografi akademik bagus.
        \item Format berkas adalah teks, sehingga mudah untuk editting
        dan dapat direkam content-tracking atau source-code manager seperti Git, CVS, atau Mercurial.
    \end{itemize}

\end{document}
